\scanpage{349}

\uline{Proposition}. Let $P$ be a section of $\mathcal M$. Equivalent conditions

\[
\begin{aligned}
\text{a)}\quad &P\text{ is regular, and }P^2=P\\
\text{b)}\quad &\det P=0,\quad\operatorname{Tr}P=1
\end{aligned}
\]

c) (If $\mathcal M=\underline{\operatorname{End}}_{\underline{\mathcal O}_S}(F)$) $\exists$ basis $e_0,e_1$ of $F$ such that the matrix of $P$ is
\[
P_{e_0,e_1}=\begin{pmatrix}0&0\\0&1\end{pmatrix}\qquad\text{i.e.}
\]
$P$ is the projection of $\underline{\mathcal O}_S e_0\oplus\underline{\mathcal O}_S e_1$ onto the factor $\underline{\mathcal O}_S e_1$

\diagram{\begin{tikzcd}[column sep=small,row sep=large,ampersand replacement=\&]
\text{d)} \& \text{\parbox{0.83\linewidth}{$\exists$ \inserted{(Zariski locally)} isomorphism $\mathcal M\xleftarrow[\sim]{\unsure{\varphi|\psi}}\mathrm M_2(\underline{\mathcal O}_S)$, taking $\begin{pmatrix}0&0\\0&1\end{pmatrix}$ to $P$.}} \\
\text{c')}\arrow[u,bend left=55] \& \text{\parbox{0.83\linewidth}{(If $\mathcal M=\underline{\operatorname{End}}_{\underline{\mathcal O}_S}(F)$) $P$ is a projector whose image \inserted{$D_1$} and kernel \inserted{$D_0$} are of rank $1$.}}
\end{tikzcd}}

\uline{Remark} \deleted{\unsure{Corollary}} For the proof, we will note that the numerical condition b) already implies that $P$ is regular, for if $P$ were a scalar $\lambda$ (on a fiber) we would have
\[
\det P=\lambda^2=0,\quad\operatorname{Tr}P=2\lambda=1
\]
hence $\lambda$ would be invertible and would have square zero, absurd. Note that this argument shows, more generally, that if $\rho$ is a section of $\mathcal M$ such that
\[
(12)\qquad \delta(\rho)=(\operatorname{Tr}\rho)^2-4\det\rho=w^2-4\unsure{v|\nu}
\]
\margin{($=$ discriminant of the Hamilton-Cayley polynomial of $\rho$)}
is an invertible section of $\underline{\mathcal O}_S$, then $\rho$ is a \uline{regular} section of $\mathcal M$.

\deleted{\uline{Corollary}} The $P$ \deleted{\unsure{in}} satisfying the conditions of the preceding proposition are called the \uline{elementary projectors} in $\mathcal M$. \unsure{Proof|their.}

\uline{Corollary} We have a one-to-one correspondence between the elementary projectors in $\mathcal M$, and the Killing pairs $(B,T)$ in $G$ (or $SB,ST$ in $SG$, or $PB,PT$ in $PG$), by

\mspage{348}

\scanpage{350}
\mspage{476}

associating to \deleted{$\unsure{N}$} $P$ the \deleted{\unsure{sub}} torus
$T=\underline{A}_P^*$, and the
\inserted{Borel $B_0$}
\deleted{unipotent subgroup}
\deleted{\inserted{additive}}
formed of the $\unsure{\alpha}$ which stabilize
\deleted{\unsure{the} flag} $D_0$
\deleted{and \unsure{induce}}
\deleted{\inserted{\unsure{the pair of the}}}
\deleted{the identity on $D_0$ and on $F/D_0\simeq D_1$,}
\inserted{hence \unsure{with respect|p.v.}}
\deleted{hence formed of those}
\deleted{\unsure{commuting} with $N$.}
to $e_0,e_1$ we will have
\[
(13)\qquad
\left\{
\begin{aligned}
T&=\left\{\begin{pmatrix}\mu_0&0\\0&\mu_1\end{pmatrix}
\ \middle|\ \mu_0,\mu_1\in\Gamma G_m\right\},\\
B_0&=\left\{\begin{pmatrix}\mu_0&\lambda\\0&\mu_1\end{pmatrix}
\ \middle|\ \mu_0,\mu_1\in\Gamma G_m,\ \lambda\in\Gamma G_a\right\}.
\end{aligned}
\right.
\]

The Killing pair \uline{opposite} to $(T,B_0)$ is
$(T,B_1)$, with
\[
(14)\qquad
B_1=\left\{\begin{pmatrix}\mu_0&0\\\lambda&\mu_1\end{pmatrix}
\ \middle|\ \mu_0,\mu_1\in\Gamma G_m,\ \lambda\in\Gamma G_a\right\},
\]
in a slightly more intrinsic way, $B_1$
\deleted{\unsure{will have}} is the stabilizer of $D_1$.
The elementary projector associated to the Killing pair
\inserted{\uline{opposite} $(T,B_1)$} is $1-N$.

\uline{Remark}. The scheme of the $\varepsilon_{\unsure{1}}$
\inserted{regular unipotents, or again of the \unsure{regular} \unsure{nilpotent} $N$}
is a bundle of punctured lines (i.e. a $G_m$-torsor)
over the scheme \inserted{$\mathcal P$} of Borels
(of $G$, of $SG$ or of $PG$, it amounts to the \unsure{same|=}),
it is \deleted{of dimension} smooth of relative dimension $2$ — the scheme of
elementary projectors $N$ is \deleted{smooth} isomorphic
to that of Killing pairs, smooth-fibered of relative dimension $1$
over the scheme of Borels. The \deleted{bundle} fiber product of
these two \inserted{relative} schemes over
$\mathcal P=$ scheme of Borels, is isomorphic to the scheme
of pinnings of $G$ (or of $SG$, or of $PG$)
i.e. of isomorphisms $\mathrm{G}\ell(2)_S\xrightarrow{\sim}G$
(resp. $\mathrm{S}\ell(2)_S\xrightarrow{\sim}SG$,

\mspage{349}

\scanpage{351}

resp. $\unsure{\mathrm{GP}(1)}_S\xrightarrow{\sim}PG$),
which is a torsor under
$\mathrm{GP}(1)_S\simeq
\underline{\mathrm{Aut}}_{S-\unsure{\mathrm{gr}}}\mathrm{G}\ell(2)_S
\simeq\underline{\mathrm{Aut}}_{S-\unsure{\mathrm{gr}}}\mathrm{S}\ell(2)_S
\simeq\underline{\mathrm{Aut}}_{S-\unsure{\mathrm{gr}}}
(\mathrm{GP}(1)_S)$.
We say that $\varepsilon$ and \deleted{$\unsure{N}$}$P$ are
\uline{associated},
\inserted{(or $N=\varepsilon-1$ and $P$ are associated),}
if they define the \unsure{same|=} Borel,
i.e. \deleted{define}
\deleted{\unsure{come}}\inserted{come} from a pinning.
We check that this amounts to the following relations
\[
(15)\qquad PN=0,\qquad NP=N
\]
\deleted{(16)} or again, in terms of $P$ and $\varepsilon=1+N$
\[
(16)\qquad P\varepsilon=P,\qquad
\varepsilon P=P+\underbrace{\varepsilon-1}_{N}
\]
\margin{This is \unsure{really} \unsure{pretty}!}

We now consider the question of
\deleted{\unsure{the equivalence}}\inserted{the} conjugation of the sections
$\widetilde\rho,\widetilde\rho'$ of $PG$,
defined by sections $\rho,\rho'$ of $G$.
We assume $\widetilde\rho$ regular i.e. $\ne1$ in
each \deleted{point} fiber, i.e. $\rho$ regular —
\deleted{\unsure{A}}
\deleted{\unsure{a}} the conjugation
\inserted{of $\widetilde\rho,\widetilde\rho'$}
by a section $u$ of $G$ (say) means that we have
\[
(17)\qquad \rho'=\lambda u\rho u^{-1},\qquad
\lambda\in\Gamma G_m
\]
whence we deduce
\[
\det\rho'=\lambda^2\det\rho\qquad
\operatorname{Tr}\rho'=\lambda\operatorname{Tr}\rho
\]
and hence
\[
(18)\qquad
\frac{\operatorname{Tr}^2\rho'}{\det\rho'}
=\frac{\operatorname{Tr}^2\rho}{\det\rho}.
\]
Let us put, for a section of $G$
\[
(19)\qquad
\deleted{\unsure{\delta|\Delta}(\rho)=}
\Delta(\rho)=\frac{\operatorname{Tr}^2\rho}{\det\rho},
\]

\mspage{350}

\scanpage{352}
\mspage{477}

we have done what was needed for \unsure{us to}
\unsure{find} a $G\longrightarrow \mathbb E^1$ which passes to the
quotient as
\[
(20)\qquad
\begin{gathered}
PG\xrightarrow{\Delta}\mathbb E^1\\
\Delta(\widetilde\rho)=\Delta(\rho)
=\frac{(\operatorname{Tr}\rho)^2}{\det\rho},
\end{gathered}
\]
and which \unsure{is|remains} invariant under inner automorphisms.

\uline{Proposition}. In order for two sections
$\widetilde\rho,\widetilde\rho'$ of $PG$ to be
\deleted{\unsure{equi}} conjugate fppf locally,
\inserted{$\widetilde\rho$ being regular,}
it is necessary \deleted{\unsure{and sufficient}}
that $\widetilde\rho'$ be \uline{regular}, and that we have
\margin{$\Big|$}
\[
(21)\qquad
\Delta(\widetilde\rho')=\Delta(\widetilde\rho)
\]
i.e. (in terms of $\rho,\rho'$ lifting
$\widetilde\rho,\widetilde\rho'$) that $\rho'$ be
regular and satisfy (18).
These conditions are sufficient when
$\Delta(\widetilde\rho)$ invertible i.e. $\operatorname{Tr}\rho$ invertible.
It remains to prove sufficiency. Let us note that fppf locally,
exists $\lambda$ \inserted{$\in\Gamma G_m$} such that
\[
\det(\rho')=\det(\lambda\rho)\qquad
\text{i.e.}\quad\lambda^2=\frac{\det\rho'}{\det\rho}.
\]
Thus \unsure{we may assume} (at the cost of replacing $\rho$ by $\lambda\rho$)
that $\det\rho=\det\rho'$.
Then (18) amounts to
\[
(22)\qquad
\operatorname{Tr}^2\rho'=\operatorname{Tr}^2\rho.
\]
\deleted{This} the question amounts to seeing that
$\exists\varepsilon\in\Gamma\mu_2$, such that
\[
\operatorname{Tr}\rho'=\varepsilon\operatorname{Tr}\rho
\quad(=\operatorname{Tr}\varepsilon\rho),
\]
we will also have
\[
\det\rho'=\det\varepsilon\rho
\quad(=\underbrace{\varepsilon^2}_{1}\det\rho=\det\rho),
\]
whence local conjugation of $\rho'$ and $\varepsilon\rho$,
we will be done. But if $\tau',\tau$ are two sections of
$\mathcal O_S$, with $\tau$ \uline{invertible}, such that
$\tau'^2=\tau^2$, we will indeed have
$\tau'=\varepsilon\tau$ $(\varepsilon=\tau'/\tau)$,
and $\varepsilon^2=1$, OK.

\mspage{351}

\scanpage{353}

\uline{Remark}. If $\Delta(\rho)$ not invertible there are
obvious counterexamples, for example with
$(\operatorname{Tr}\rho)^2=0$,
$(\operatorname{Tr}\rho')^2=0$,
$\det\rho=\det\rho'=1$,
$\operatorname{Tr}\rho$ and $\operatorname{Tr}\rho'$
not being homothetic. We can take
\[
\rho=\begin{pmatrix}0&1\\-1&w\end{pmatrix},
\qquad
\rho'=\begin{pmatrix}0&1\\-1&w'\end{pmatrix}
\]
$w^2=w'^2=0$, $w,w'$ not homothetic, for example
base ring
$A=A_0[w,w']/(w^2,ww',w'^2)$
(dual \unsure{numbers|numbers} on the free module of rank $2$
generated by $w,w'$).

$2)$ Relative position of two sections $\rho,\sigma$
of $G$, or $\widetilde\rho,\widetilde\sigma$ the two sections of $PG$.
\margin{$\big|$}

In $PG$ we are interested in
\[
(23)\qquad
\widetilde\varepsilon_0
=\widetilde\sigma^{-1}\widetilde\rho
\quad\text{and}\quad
\widetilde\varepsilon_1
=\widetilde\rho\widetilde\sigma^{-1}
=\widetilde\rho(\widetilde\varepsilon_0)
=\widetilde\sigma(\widetilde\varepsilon_1)
\]
and in the condition that
\inserted{\textcircled{\(1^\circ\)}}
$\widetilde\varepsilon_0$,
\deleted{\unsure{$\widetilde\varepsilon_1$, both}}
\inserted{(hence $\widetilde\varepsilon_1$)}
\unsure{be|are} regular unipotent —
hence \inserted{$\widetilde\varepsilon_0,\widetilde\varepsilon_1$}
come canonically from the
regular unipotents $\varepsilon_0,\varepsilon_1$ of $SG$,
satisfying
\[
(24)\qquad
\varepsilon_1=\widetilde\rho(\varepsilon_0)
=\widetilde\sigma(\varepsilon_0)
\]
\[
(25)\qquad
\left\{
\begin{aligned}
\varepsilon_0&=(\text{Invertible scalar})\,\sigma^{-1}\rho\\
\varepsilon_1&=(-\ -\ -\ -)\,\rho\sigma^{-1}
\end{aligned}
\right.
\]
\margin{for a \unsure{more}
\unsure{convenient|suitable} \unsure{lifting} of $\sigma$ and $\rho$}
\margin{i.e.
$\varepsilon_0=-\deleted{\unsure{\lambda_0}}\sigma^{-1}\rho$,
$\varepsilon_1=-\deleted{\unsure{\lambda_0}}\rho\sigma^{-1}$}
and \textcircled{\(2^\circ\)}
$\widetilde\varepsilon_0,\widetilde\varepsilon_1$ be
\deleted{everywhere}\inserted{in each fiber}
``non-equivalent'',
i.e. \deleted{$\varepsilon_0=$}, let us put
\[
(26)\qquad
\varepsilon_0=1+N_0,\qquad
\varepsilon_1=1+N_1
\]

\mspage{352}